\documentclass[11pt,a4paper]{article}

\usepackage[margin=1in]{geometry}
\usepackage{amsmath,amssymb}
\usepackage{booktabs}
\usepackage{listings}
\usepackage{xcolor}
\usepackage{hyperref}

\definecolor{codegray}{gray}{0.95}
\lstset{
  basicstyle=\ttfamily\small,
  backgroundcolor=\color{codegray},
  frame=single,
  breaklines=true,
  columns=fullflexible
}

\title{Assignment 1: MicroGrad\\Scalar Automatic Differentiation Engine}
\author{John George Mikhael \and 23P0266}
\date{\today 3/11/2026}

\begin{document}
\maketitle

\begin{abstract}
This report describes a scalar automatic differentiation engine and a small
multilayer perceptron implemented without PyTorch or another machine-learning
framework. The implementation represents each scalar as a node in a directed
computational graph, records local derivative rules for elementary operations,
and computes gradients by reverse-mode automatic differentiation.
\end{abstract}

\section{Objective}
The objective of this assignment is to understand the mechanics behind
backpropagation. The implementation in \texttt{src/engine.py} provides a
\texttt{Value} class for scalar data and gradients. The implementation in
\texttt{src/nn.py} builds a simple neural-network API consisting of
\texttt{Neuron}, \texttt{Layer}, and \texttt{MLP}.

\section{Computational Graph}
Each \texttt{Value} stores:
\begin{itemize}
  \item \texttt{data}: the scalar value of the node;
  \item \texttt{grad}: the derivative of the final output with respect to the
        node;
  \item \texttt{\_prev}: the set of parent nodes used to compute it;
  \item \texttt{\_backward}: a closure containing the local chain-rule update.
\end{itemize}

For example, if $z=x+y$, then
\[
\frac{\partial z}{\partial x}=1,\qquad
\frac{\partial z}{\partial y}=1.
\]
When the output gradient is $\partial L/\partial z$, the local backward
function accumulates
\[
\frac{\partial L}{\partial x}
\mathrel{+}=\frac{\partial L}{\partial z},\qquad
\frac{\partial L}{\partial y}
\mathrel{+}=\frac{\partial L}{\partial z}.
\]
The use of $\mathrel{+}=$ is essential when a value is reused in multiple
branches of the graph.

\section{Local Derivative Rules}
The engine implements the following rules:
\begin{align*}
z &= x+y &
\frac{\partial z}{\partial x} &= 1,\quad
\frac{\partial z}{\partial y}=1,\\
z &= xy &
\frac{\partial z}{\partial x} &= y,\quad
\frac{\partial z}{\partial y}=x,\\
z &= x^p &
\frac{\partial z}{\partial x} &= p x^{p-1},\\
z &= \tanh(x) &
\frac{\partial z}{\partial x} &= 1-\tanh^2(x),\\
z &= e^x &
\frac{\partial z}{\partial x} &= e^x.
\end{align*}
Subtraction, negation, and division are reduced to these primitives:
\[
x-y=x+(-y),\qquad
-x=(-1)x,\qquad
\frac{x}{y}=x y^{-1}.
\]

\section{Topological Sort and Backpropagation}
The graph is traversed by depth-first search. A visited set ensures that a
shared node is visited once, while post-order insertion ensures that every
parent appears before its dependent node in the resulting topological list.
For a scalar output $L$, the global backward algorithm is:
\begin{enumerate}
  \item Build the topological ordering starting at $L$.
  \item Set $L.\texttt{grad}=1$ because $\partial L/\partial L=1$.
  \item Traverse the ordering in reverse and invoke each node's local
        \texttt{\_backward} closure.
\end{enumerate}
This order guarantees that a node receives all gradient contributions from
its children before its own gradient is propagated to its parents.

\subsection{Diamond graph}
Consider $c=a+a$ and $d=ca$. Since $d=2a^2$,
\[
\frac{\partial d}{\partial a}=4a.
\]
At $a=3$, the correct result is $12$. The implementation obtains this result
because both paths from $a$ are accumulated rather than overwritten.

\section{Neural-Network API}
A \texttt{Neuron} contains a list of randomly initialized weights and one
randomly initialized bias. Given an input vector $\boldsymbol{x}$, it computes
\[
y=\tanh\left(\sum_i w_i x_i+b\right).
\]
A \texttt{Layer} contains multiple neurons and returns either a single
\texttt{Value} for one output neuron or a list of values for multiple output
neurons. An \texttt{MLP} creates consecutive layers using the requested size
configuration, such as \texttt{MLP(2,[4,4,1])}.

The \texttt{parameters()} methods flatten all weights and biases so that a
training loop can update them uniformly. The \texttt{zero\_grad()} methods
clear stored gradients before a new optimization step.

\section{Training Demonstration}
The accompanying notebook, \texttt{demo\_micrograd.ipynb}, trains
\texttt{MLP(2,[4,4,1])} on the four XOR examples:
\[
(0,0)\mapsto -1,\quad
(0,1)\mapsto +1,\quad
(1,0)\mapsto +1,\quad
(1,1)\mapsto -1.
\]
The output neuron uses $\tanh$, so the targets lie in $[-1,1]$. The mean
squared error is
\[
\mathcal{L}=\frac{1}{N}\sum_{j=1}^{N}(\hat{y}_j-y_j)^2.
\]
Each training step resets gradients, evaluates the loss, calls
\texttt{backward()}, and applies manual gradient descent:
\[
\theta \leftarrow \theta-\eta\nabla_\theta\mathcal{L}.
\]
The notebook uses a fixed random seed for reproducibility and plots both the
loss curve and the learned XOR predictions.

\section{Testing and Results}
The provided automated test suite contains 12 tests. It checks initialization,
arithmetic operations, invalid operand handling, powers, activations,
gradient accumulation, chain-rule behavior, division by zero, type casting,
and the complete neural-network API. All 12 tests pass after implementation.

\begin{table}[h]
\centering
\caption{Implemented components and validation}
\begin{tabular}{ll}
\toprule
Component & Result\\
\midrule
Scalar arithmetic and reverse operators & Implemented\\
Power, \texttt{tanh}, and \texttt{exp} & Implemented\\
Topological reverse-mode differentiation & Implemented\\
Neuron, Layer, and MLP & Implemented\\
Automated tests & 12 passed\\
\bottomrule
\end{tabular}
\end{table}

\section{Design Justifications}
\begin{itemize}
  \item Scalars are represented as graph nodes so each operation can retain
        both its forward result and local derivative rule.
  \item Parent nodes are stored in a set because graph membership is what is
        needed for traversal; gradient contributions are accumulated inside
        the backward closures.
  \item Reverse-mode differentiation is appropriate because the network has
        one scalar loss and many parameters.
  \item The neural-network API intentionally stays small: parameter
        collection and gradient reset are sufficient for a transparent
        hand-written optimization loop.
\end{itemize}

\section{Conclusion}
This assignment demonstrates that backpropagation can be implemented with
basic Python objects, local derivative formulas, and a topological traversal.
The same chain rule used by this small engine is the foundation of gradient
computation in larger deep-learning frameworks.

\end{document}
